\documentclass[a4paper,11pt]{article}
\usepackage[utf8]{inputenc}
\usepackage[T1]{fontenc}
\usepackage[english]{babel}
\usepackage[left=1cm,right=1cm,top=.5cm,bottom=0cm]{geometry}
\usepackage{enumitem}
\usepackage{hyperref}
\usepackage{xcolor}
\usepackage{titlesec}
\usepackage{marvosym}
\usepackage{lmodern}
\usepackage[normalem]{ulem}

% --- Couleurs ---
\definecolor{primary}{RGB}{0, 45, 114}
\definecolor{secondary}{RGB}{80, 80, 80}

% --- Config Liens ---
\hypersetup{colorlinks=true, linkcolor=primary, filecolor=primary, urlcolor=primary}

% --- Titres ---
\titleformat{\section}{\large\bfseries\color{primary}\uppercase}{}{0em}{}[\titlerule]
\titlespacing{\section}{0pt}{8pt}{5pt}

\begin{document}

% --- EN-TÊTE ---
\begin{center}
    {\Large \textbf{Lo\"ic Aron MBASSI EWOLO}} \\ \vspace{4pt}
    \textbf{\large Concepteur d'Agents IA pour l'Amelioration Continue} \\
    \textit{\small Stage INRIA Saclay (Avril--Ao\^ut 2026) -- Available for 12-month apprenticeship from September 2026} \\ \vspace{4pt}
    \small
    \Pointinghand\ Cergy, France (Mobile IDF) \quad
    \Mobilefone\ (+33) 7 59 52 33 98 \quad
    {\color{primary}\Letter\ \href{mailto:loicaron.mbassiewolo@ensea.fr}{loicaron.mbassiewolo@ensea.fr}} \\ \vspace{2pt}
    {\color{primary}LinkedIn: \href{https://www.linkedin.com/in/loic-mbassi/}{linkedin.com/in/loic-mbassi}} \quad
    {\color{primary}GitHub: \href{https://github.com/Nameless0l}{github.com/Nameless0l}} \quad
    {\color{primary}Portfolio: \href{https://mbassiloic.tech}{mbassiloic.tech}}
\end{center}

\vspace{-6pt}

% --- PROFIL ---
\noindent\small{Designing intelligent agents for continuous improvement: RAG-powered knowledge systems, multi-agent orchestration with Google ADK, NLP analysis, optimization algorithms for resource efficiency.}

% --- FORMATION ---
\section{Formation}
\begin{itemize}[leftmargin=*, label=\tiny$\bullet$, nosep]
    \item \textbf{ENSEA -- Ecole Nationale Sup\'erieure de l'Electronique et de ses Applications} \hfill \textit{2025 -- 2027} \\
    \small{Double diplome d'Ingenieur -- Syst\`emes Embarqu\'es, Signal, IA.}
    \item \textbf{Ecole Polytechnique de Yaounde\'} \hfill \textit{2021 -- 2025} \\
    \small{Diplome d'Ingenieur de Conception: Mathematiques Appliquees, Genie Logiciel.}
\end{itemize}

% --- COMPETENCES ---
\section{Competences Techniques}
\begin{itemize}[leftmargin=*, nosep]
    \item \small{\textbf{AI \& Agents:} Python, LangChain, RAG, Gemini, Google ADK, NLP, orchestration}
    \item \small{\textbf{Backend \& Data:} FastAPI, Docker, SQL, Git, optimization algorithms, scalability}
\end{itemize}

% --- EXPERIENCES ---
\section{Projets / Experiences}
\begin{itemize}[leftmargin=*, label={-}]

\item \textbf{Pipeline NLP et Analyse Documentaire a Grande Echelle} \hfill \textit{Avr 2026 -- Ao\^ut 2026} \\
\textit{INRIA Saclay -- Equipe TRIBE} \hfill \textit{Python, NLP, BERT, pandas, Git}
\vspace{-2pt}
    \begin{itemize}[leftmargin=0.4cm, nosep, label=\tiny$\bullet$]
        \item \small{Conception d'une infrastructure NLP pour analyse de documents complexes a grande echelle.}
        \item \small{Scoring interpretable agr\'egeant analyses multiples pour decision support systems.}
    \end{itemize}
\vspace{3pt}

\item \textbf{Systeme Multi-Agents Collaboratif avec RAG} \hfill \textit{2024 -- 2025} \\
\textit{Recherche Google ADK} \hfill \textit{Google ADK, LangChain, RAG, Gemini, Docker}
\vspace{-2pt}
    \begin{itemize}[leftmargin=0.4cm, nosep, label=\tiny$\bullet$]
        \item \small{Orchestration de 5 agents IA specialises avec resolution dynamique de conflits.}
        \item \small{Base de connaissances enrichie avec RAG; recommendations multi-criteres optimisees.}
    \end{itemize}
\vspace{3pt}

\item \textbf{Architecture Metier avec Domain-Driven Design} \hfill \textit{2024 -- 2025} \\
\textit{Snappy -- Enterprise Messaging Platform} \hfill \textit{Spring Boot, Domain-Driven Design, PostgreSQL, Redis}
\vspace{-2pt}
    \begin{itemize}[leftmargin=0.4cm, nosep, label=\tiny$\bullet$]
        \item \small{Conception d'une architecture DDD scalable avec agrégats de domaine, invariants, machine a etats.}
        \item \small{Chatbots IA avec RAG pour reponses contextuelles basees sur documents metier.}
    \end{itemize}
\vspace{3pt}

\item \textbf{Optimisation ML pour Collecte et Logistique} \hfill \textit{2024} \\
\textit{1er Prix CITS National Hackathon} \hfill \textit{Python, ML, optimization algorithms, IoT}
\vspace{-2pt}
    \begin{itemize}[leftmargin=0.4cm, nosep, label=\tiny$\bullet$]
        \item \small{Algorithmes d'optimisation (TSP) pour reduction des coûts logistiques de 20\%.}
        \item \small{Modeles predictifs ML pour prevision de demande et planification intelligente.}
    \end{itemize}

\end{itemize}

\end{document}
